% =====================================================================
% Capítulo 3 — Trabalhos Relacionados
% Dissertação de Mestrado MAS-Hunt (ABNT, PT-BR)
% Convenções: ver .orchestration/dissertacao-OUTLINE.md
% Citações: chaves reais de bibliografia.bib e mashunt/references.bib
% NOTA DE INTEGRAÇÃO (cite gap): chaves dos eixos (a),(b),(d) vêm de
%   bibliografia.bib; chaves do eixo (c) — adversariais — vêm de
%   mashunt/references.bib. As duas bases são disjuntas; o main.tex deve
%   apontar \bibliography para ambas (ou mesclá-las) para resolver todas as
%   chaves abaixo. A tabela usa \ding{} do pacote pifont.
% =====================================================================

\chapter{Trabalhos Relacionados}
\label{cap:trabalhos-relacionados}

A proposta do MAS-Hunt situa-se na confluência de quatro linhas de pesquisa que,
até o momento, evoluíram de modo amplamente independente: (i) a aplicação de
\textit{Large Language Models} (LLMs) e agentes autônomos à caça a ameaças e à
detecção comportamental; (ii) os sistemas multiagente para segurança
cibernética e sua governança; (iii) os ataques adversariais dirigidos aos
próprios agentes baseados em LLM, em especial o envenenamento de memória e a
injeção de prompt/logs; e (iv) a automação de detecção em ambientes
SIEM/SOC. Este capítulo organiza a literatura segundo esses quatro eixos
temáticos, compara criticamente as abordagens dentro de cada eixo e, ao final,
explicita a lacuna que o MAS-Hunt se propõe a preencher. A revisão é deliberadamente
comparativa: o objetivo não é catalogar trabalhos, mas evidenciar que nenhuma das
linhas existentes integra, em um único artefato avaliado empiricamente,
governança multiagente real e resistência a injeção adversarial sobre telemetria
viva.

% =====================================================================
\section{Caça a Ameaças e Detecção Comportamental Assistidas por IA}
\label{sec:rel-hunting-ia}

A caça a ameaças (\textit{threat hunting}) consolidou-se como um deslocamento
paradigmático da postura reativa, baseada em assinaturas, para uma busca
proativa e iterativa por adversários já presentes no ambiente
\cite{hillier2022turninghuntedhunterthreat}. \citeonline{hillier2022turninghuntedhunterthreat}
descrevem o ciclo de vida da caça a ameaças e seu ecossistema, antecipando a
\textit{grande promessa} da inteligência artificial como força multiplicadora do
analista humano; o trabalho, contudo, permanece no plano conceitual, sem
implementação operacional. Essa promessa é sistematizada de forma prospectiva por
\citeonline{academia22}, que propõe a \textit{caça autônoma a ameaças} como
paradigma futuro de inteligência de ameaças dirigida por IA, e por
\citeonline{academia21}, que mapeia o \textit{caminho para a defesa cibernética
autônoma}. Ambos os trabalhos, entretanto, são essencialmente revisões de visão:
identificam direções desejáveis sem entregar uma arquitetura validada nem
quantificar resiliência.

No plano das técnicas concretas de detecção, há uma trajetória clara que vai do
aprendizado de máquina supervisionado em domínio fechado até o uso de LLMs como
motores de raciocínio. \citeonline{HADDADPAJOUH201888} aplicam redes neurais
recorrentes profundas à caça a \textit{malware} em dispositivos de Internet das
Coisas, alcançando bons índices de classificação, porém à custa de um modelo
treinado de modo \textit{offline}, incapaz de formular hipóteses ou de interagir
com a telemetria em tempo de execução. As limitações intrínsecas dessa família de
abordagens são sintetizadas por \citeonline{ucci2019}, cujo levantamento sobre
técnicas de aprendizado de máquina para análise de \textit{malware} evidencia a
dependência de \textit{features} predominantemente sintáticas e a fragilidade
frente à mudança de distribuição entre o ambiente de treino e o \textit{endpoint}
real.

A introdução de LLMs altera qualitativamente esse panorama ao acoplar raciocínio
em linguagem natural e capacidade de geração de consultas. O HuntGPT
\cite{ali2023huntgptintegratingmachinelearningbased} é representativo dessa
geração: integra detecção de anomalias baseada em aprendizado de máquina,
inteligência artificial explicável e um LLM como camada de interpretação,
oferecendo ao analista justificativas legíveis para os alertas. Trata-se,
porém, de um agente único e assistivo, sem orquestração colaborativa nem
qualquer hipótese de defesa contra manipulação do próprio modelo. O modelo de
geração de hipóteses de \citeonline{10713082} formaliza o estágio mais cognitivo
da caça --- a formulação de hipóteses verificáveis ---, mas o faz como
componente isolado, desacoplado de execução, validação ou resposta. O TIRA, de
\citeonline{10750987}, avança rumo à automação de ponta a ponta (caça, detecção e
perfilamento de atores), aproximando-se de um fluxo operacional completo;
todavia, não modela um adversário inteligente capaz de subverter o próprio
sistema de caça, premissa central do MAS-Hunt.

Em síntese, o eixo de caça assistida por IA progrediu da classificação
\textit{offline} \cite{HADDADPAJOUH201888,ucci2019} para a interpretação
explicável \cite{ali2023huntgptintegratingmachinelearningbased} e a geração de
hipóteses \cite{10713082,10750987}, mas mantém duas omissões recorrentes: a
ausência de \textit{colaboração estruturada entre múltiplos agentes
especializados} e a desconsideração da \textit{superfície de ataque do próprio
sistema de IA}.

% =====================================================================
\section{Sistemas Multiagente para Segurança e sua Governança}
\label{sec:rel-mas-governanca}

A coordenação entre múltiplos agentes para resolver problemas distribuídos é um
campo maduro, cujas raízes remontam à formalização de organizações, planos e
escalonamento como mecanismos de coordenação de agentes de IA
\cite{durfee1993organizations}. O advento dos LLMs reabilitou esse campo sob nova
forma: agentes de linguagem passaram a coordenar-se em arranjos hierárquicos para
tarefas complexas. \citeonline{liu2023llm} propõem um agente hierárquico de
linguagem para coordenação humano-IA em tempo real, demonstrando que a
decomposição em camadas --- um agente de alto nível que delega a subagentes ---
melhora a latência e a qualidade da cooperação. \citeonline{agashe2023llmcoordination}
avaliam sistematicamente as capacidades de coordenação multiagente de diferentes
LLMs, e \citeonline{qiu2024collaborativeintelligencepropagatingintentions}
mostram que a propagação explícita de intenções e raciocínio entre agentes
aprimora a inteligência coletiva. Padrões cooperativos correlatos aparecem no
ProAgent \cite{DBLP:journals/corr/abs-2308-11339}, na coordenação \textit{zero-shot}
de \citeonline{pmlr-v202-li23au} e na coordenação fundamentada em modelos
multimodais de \citeonline{mahmud2025distributedmultiagentcoordinationusing}.
Essa literatura, contudo, otimiza \textit{desempenho cooperativo} em ambientes
benignos: a coordenação é tratada como problema de eficiência
\cite{bai2024reducingredundantcomputationmultiagent}, não de segurança.

Um subconjunto da literatura começa a tratar a coordenação como problema
adversarial e de confiança. \citeonline{Yuan_Zhang_Xue_Yin_Chen_Guan_Li_Qian_Yu_2023}
propõem coordenação multiagente robusta via geração evolutiva de atacantes
auxiliares, endurecendo políticas contra perturbações; o BlockAgents
\cite{doe2023emergent} busca coordenação tolerante a falhas bizantinas por meio de
\textit{blockchain}, atacando o problema do agente comprometido pela via da
imutabilidade do registro. Esses trabalhos, entretanto, situam-se no domínio de
aprendizado por reforço multiagente ou de consenso distribuído genérico, sem
ancoragem em telemetria de segurança real nem nos vetores de ataque específicos a
LLMs (envenenamento de memória, injeção de prompt).

Paralelamente, emerge uma literatura de \textit{governança} de sistemas de IA que
fornece o vocabulário normativo --- mas não a implementação técnica --- para o
controle de agentes autônomos. \citeonline{isaca2025leveraging} adaptam o
arcabouço COBIT à governança de sistemas de IA; \citeonline{mcintosh2024cobit} e
\citeonline{McIntosh2024} avaliam a transição do COBIT para a ISO/IEC 42001 sob a
ótica de oportunidades, riscos e conformidade regulatória na comercialização de
LLMs. Esses estudos estabelecem \textit{por que} a governança é necessária e
\textit{quais} controles deveriam existir, mas não descem ao nível de protocolos
executáveis de mensagens tipadas, votação de conselho ou quarentena de agentes ---
exatamente o terreno onde o MAS-Hunt opera. A obra de referência sobre
\textit{frameworks} de agentes \cite{guo2024survey} complementa esse quadro no
plano da engenharia, mas sem viés de segurança.

A comparação entre os dois subcampos revela uma assimetria reveladora: a
coordenação multiagente sabe \textit{como} fazer agentes cooperarem, e a
governança de IA sabe \textit{quais} controles são desejáveis, porém raramente os
dois se encontram em um sistema concreto de segurança cibernética cuja
arquitetura de governança seja ela própria o mecanismo de defesa.

% =====================================================================
\section{Ataques Adversariais a Agentes Baseados em LLM}
\label{sec:rel-ataques-adversariais}

O terceiro eixo é o que motiva diretamente o endurecimento do MAS-Hunt: a
constatação empírica de que a autonomia que torna os agentes poderosos os torna
igualmente frágeis. A superfície de ataque mais crítica é a memória cognitiva do
agente. \citeonline{NEURIPS2024_eb113910}, com o AgentPoison, demonstram que a
memória ou a base de conhecimento de recuperação aumentada (RAG) de um agente
pode ser envenenada com taxa de injeção inferior a 0{,}1\% para induzir
comportamento malicioso direcionado com taxa de sucesso de ataque superior a
80\%. Esse resultado é decisivo: mostra que a integridade da memória não pode ser
presumida e fundamenta o mecanismo M1 (Integridade de Memória) do MAS-Hunt.

Uma segunda classe de ataques explora o \textit{loop de raciocínio} do agente.
\citeonline{zhang2024breaking}, em \textit{Breaking Agents}, caracterizam a
\textit{amplificação de mau funcionamento} (\textit{malfunction amplification}),
em que um erro menor desencadeia um laço de realimentação de ações inúteis, com
taxas de falha que ultrapassam 80\% --- um vetor efetivo de negação de serviço
contra o próprio agente. A injeção indireta de \textit{prompt}, por sua vez, é
quantificada pelo \textit{benchmark} InjecAgent \cite{zhan2024injecagent}, que
mede a vulnerabilidade de agentes integrados a ferramentas a instruções
maliciosas embutidas em conteúdo externo --- precisamente o cenário da
injeção adversarial de logs, em que o adversário contamina a telemetria que o
caçador irá ler. O CRDA \cite{liu2024crda} estende a análise ao \textit{drift} de
risco de conteúdo via interação adversarial multiagente, evidenciando que a
degradação pode emergir da própria dinâmica entre agentes.

O caráter sistêmico dessas ameaças em \textit{coletivos} de agentes é o achado
mais alarmante. \citeonline{gu2024agent}, com o Agent Smith, mostram que uma
única imagem adversarial pode desencadear um \textit{jailbreak infeccioso} que se
propaga exponencialmente por um milhão de agentes multimodais.
\citeonline{huang2024resilience} estudam a resiliência de sistemas multiagente na
presença de agentes maliciosos e apresentam o resultado de projeto mais relevante
para esta dissertação: estruturas de orquestração \textit{hierárquicas} exibem a
menor degradação de desempenho sob condições adversariais quando comparadas a
topologias alternativas. Esse achado justifica diretamente a escolha arquitetural
do MAS-Hunt por uma hierarquia de três camadas. Por fim, essas vulnerabilidades
devem ser lidas à luz de um adversário ativo e inteligente:
\citeonline{moreno2025analysis} propõem teste de intrusão autônomo via
aprendizado por reforço e sistemas de recomendação, e \citeonline{academia20}
discute a IA como o \textit{novo atacante}, desenvolvendo agentes para segurança
ofensiva. Um agente ofensivo dessa natureza é o candidato natural a executar os
ataques de envenenamento e de manipulação comportamental acima descritos --- donde
a conclusão de que qualquer arquitetura defensiva que não modele um adversário de
IA igualmente capaz é fundamentalmente incompleta.

A leitura conjunta deste eixo estabelece a tríade de ameaças que o MAS-Hunt
endereça: envenenamento de memória \cite{NEURIPS2024_eb113910}, amplificação de
mau funcionamento \cite{zhang2024breaking} e injeção/propagação infecciosa
\cite{zhan2024injecagent,gu2024agent,huang2024resilience}. Crucialmente, a
literatura deste eixo \textit{ataca} e \textit{mede}, mas raramente \textit{defende}:
as defesas práticas, quando propostas, permanecem em ambiente simulado, desacopladas
de qualquer tarefa de segurança operacional como a caça a ameaças.

% =====================================================================
\section{Automação de Detecção em Ambientes SIEM/SOC}
\label{sec:rel-siem-soc}

O quarto eixo trata da automação da detecção no contexto operacional de um centro
de operações de segurança (SOC) apoiado por um \textit{Security Information and
Event Management} (SIEM). A caça a \textit{ransomware} via técnicas de
\textit{threat hunting} é levantada por \citeonline{9791234}, que sistematiza o
estado da arte e expõe a dependência de regras e heurísticas construídas
manualmente. A análise de logs como substrato primário da caça é tratada por
\citeonline{10616474}, evidenciando tanto o valor da telemetria bruta quanto sua
suscetibilidade à manipulação --- um log adulterado é um vetor de injeção
adversarial direto. Em ambientes distribuídos, \citeonline{9931222} aplicam
aprendizado federado à caça a ameaças para detecção de ataques de Ameaça
Persistente Avançada (APT) em redes definidas por software, demonstrando
detecção colaborativa que preserva privacidade, porém sem agentes de raciocínio
nem governança sobre eles.

O obstáculo estatístico que assombra toda automação de detecção em larga escala é
a falácia da taxa de base bayesiana, articulada por \citeonline{dolan-gavitt2025}:
mesmo modelos altamente acurados geram um volume avassalador de falsos positivos
quando o evento procurado é raro, tornando-os operacionalmente inúteis. A resposta
proposta naquele trabalho --- validação determinística baseada em \textit{canaries}
e provas de posse, que mira uma taxa de falsos positivos próxima de zero ---
informa diretamente o mecanismo M2 (Validação Cruzada) do MAS-Hunt, que exige
corroboração por evidência verificável antes de qualquer escalonamento. No plano
da detecção comportamental sobre o \textit{endpoint},
A literatura de análise de malware documenta a sensibilidade de detectores a
traços comportamentais e ao contexto de execução \cite{ucci2019}, o que motiva a
opção do MAS-Hunt por operar diretamente sobre telemetria viva da pilha Elastic,
e não apenas sobre traços sintéticos. Abordagens correlatas que integram LLMs a
detecção anômala explicável \cite{ali2023huntgptintegratingmachinelearningbased}
reforçam a tendência, mas mantêm o paradigma de modelo único, sem orquestração
nem \textit{hardening}.

A comparação neste eixo expõe um padrão consistente: a automação SIEM/SOC ou
delega a regras determinísticas frágeis e ruidosas \cite{9791234,10616474}, ou
emprega modelos estatísticos vulneráveis à falácia da taxa de base
\cite{dolan-gavitt2025} e à mudança de distribuição \cite{ucci2019}. Em
nenhum dos casos a própria camada de automação é tratada como alvo a ser
endurecido contra manipulação adversarial.

% =====================================================================
\section{Síntese Comparativa}
\label{sec:rel-sintese}

A Tabela~\ref{tab:rel-comparacao} confronta representantes de cada eixo segundo
cinco dimensões que delimitam a contribuição do MAS-Hunt: orquestração multiagente
real (não um agente único nem um \textit{pipeline} monolítico); governança
explícita (conselho, mensagens tipadas, papéis); resiliência a ataques
adversariais ao próprio agente (memória, injeção, mau funcionamento); operação
sobre telemetria viva de produção; e avaliação empírica controlada da resiliência.

\begin{table}[htbp]
\centering
\caption{Comparação dos trabalhos relacionados segundo as dimensões críticas do MAS-Hunt. Legenda: \ding{51}~contemplado; $\sim$~parcial/conceitual; \ding{55}~não contemplado.}
\label{tab:rel-comparacao}
\small
\begin{tabular}{p{4.4cm}ccccc}
\hline
\textbf{Trabalho / Eixo} & \textbf{Multi-} & \textbf{Gover-} & \textbf{Resil.} & \textbf{Telem.} & \textbf{Aval.} \\
                         & \textbf{agente} & \textbf{nança}  & \textbf{advers.} & \textbf{viva}  & \textbf{empír.} \\
\hline
\multicolumn{6}{l}{\textit{(a) Caça/detecção assistida por IA}} \\
HuntGPT \cite{ali2023huntgptintegratingmachinelearningbased}        & \ding{55} & \ding{55} & \ding{55} & $\sim$    & \ding{51} \\
Ger.\ de hipóteses \cite{10713082}                                  & \ding{55} & \ding{55} & \ding{55} & \ding{55} & $\sim$    \\
TIRA \cite{10750987}                                                & $\sim$    & \ding{55} & \ding{55} & $\sim$    & \ding{51} \\
Caça autônoma \cite{academia22}                                     & $\sim$    & \ding{55} & \ding{55} & \ding{55} & \ding{55} \\
\hline
\multicolumn{6}{l}{\textit{(b) Multiagente e governança}} \\
Agente hierárquico \cite{liu2023llm}                                & \ding{51} & $\sim$    & \ding{55} & \ding{55} & \ding{51} \\
BlockAgents \cite{doe2023emergent}                                  & \ding{51} & $\sim$    & $\sim$    & \ding{55} & \ding{51} \\
Coord.\ robusta \cite{Yuan_Zhang_Xue_Yin_Chen_Guan_Li_Qian_Yu_2023} & \ding{51} & \ding{55} & $\sim$    & \ding{55} & \ding{51} \\
Governança COBIT/ISO \cite{mcintosh2024cobit}                       & \ding{55} & \ding{51} & \ding{55} & \ding{55} & \ding{55} \\
\hline
\multicolumn{6}{l}{\textit{(c) Ataques adversariais a agentes}} \\
AgentPoison \cite{NEURIPS2024_eb113910}                             & $\sim$    & \ding{55} & \ding{55} & \ding{55} & \ding{51} \\
Breaking Agents \cite{zhang2024breaking}                            & $\sim$    & \ding{55} & \ding{55} & \ding{55} & \ding{51} \\
Agent Smith \cite{gu2024agent}                                      & \ding{51} & \ding{55} & \ding{55} & \ding{55} & \ding{51} \\
Resiliência MAS \cite{huang2024resilience}                          & \ding{51} & $\sim$    & $\sim$    & \ding{55} & \ding{51} \\
\hline
\multicolumn{6}{l}{\textit{(d) Automação SIEM/SOC}} \\
Caça a \textit{ransomware} \cite{9791234}                           & \ding{55} & \ding{55} & \ding{55} & $\sim$    & \ding{55} \\
Caça federada APT \cite{9931222}                                    & \ding{51} & \ding{55} & \ding{55} & $\sim$    & \ding{51} \\
Validação determin.\ \cite{dolan-gavitt2025}                        & \ding{55} & \ding{55} & $\sim$    & $\sim$    & \ding{51} \\
\hline
\textbf{MAS-Hunt (este trabalho)} & \ding{51} & \ding{51} & \ding{51} & \ding{51} & \ding{51} \\
\hline
\end{tabular}
\end{table}

A leitura da Tabela~\ref{tab:rel-comparacao} torna a lacuna manifesta. Os trabalhos
do eixo (a) operam, na melhor das hipóteses, com um agente único sobre telemetria,
sem governança nem modelo de adversário. Os do eixo (b) dominam a orquestração e o
discurso de governança, mas não os confrontam com ataques concretos a agentes nem
com telemetria de segurança real. Os do eixo (c) demonstram exaustivamente \textit{que}
os agentes são vulneráveis, porém em laboratório simulado e sem propor defesa
operacional integrada a uma tarefa de caça. Os do eixo (d) automatizam detecção,
mas tratam o componente de IA como caixa estatística, jamais como alvo a ser
endurecido. Nenhuma linha, isoladamente, ocupa simultaneamente as cinco colunas.

% =====================================================================
\section{A Lacuna que o MAS-Hunt Preenche}
\label{sec:rel-lacunas}

Da análise comparativa precedente decorre a lacuna central que esta dissertação
endereça. As quatro linhas de pesquisa convergem para um mesmo ponto cego: cada
uma resolve uma fração do problema, mas nenhuma integra, em um único sistema
avaliado empiricamente, \textit{governança multiagente real} e \textit{resistência
a injeção adversarial} sobre dados de produção. Especificamente, identificam-se
três deficiências que o MAS-Hunt se propõe a sanar:

\begin{enumerate}
    \item \textbf{Ausência de orquestração multiagente com governança como mecanismo
    de defesa.} A literatura de coordenação \cite{liu2023llm,agashe2023llmcoordination,qiu2024collaborativeintelligencepropagatingintentions}
    otimiza cooperação, e a de governança \cite{isaca2025leveraging,mcintosh2024cobit}
    prescreve controles, mas não há sistema em que a hierarquia de três camadas
    --- Conselho de Governança $\rightarrow$ Gerentes $\rightarrow$ Caçadores ---,
    com mensagens tipadas e protocolos de consenso, atue \textit{ela própria} como
    a barreira que impede a propagação de um agente comprometido. O achado de
    \citeonline{huang2024resilience} de que estruturas hierárquicas degradam menos
    sob ataque é aqui transformado de observação em princípio de projeto operante.

    \item \textbf{Falta de validação de integridade em tempo de execução contra
    ataques a agentes.} Os trabalhos que demonstram envenenamento de memória
    \cite{NEURIPS2024_eb113910}, injeção indireta \cite{zhan2024injecagent} e
    amplificação de mau funcionamento \cite{zhang2024breaking} caracterizam as
    ameaças sem oferecer contramedidas operacionais acopladas a uma tarefa real. O
    MAS-Hunt responde com os mecanismos M1 (Integridade de Memória), M2 (Validação
    Cruzada por evidência verificável, na linha da validação determinística de
    \citeonline{dolan-gavitt2025}), M3 (\textit{framework} de análise estruturada
    OBSERVAR $\rightarrow$ HIPOTETIZAR $\rightarrow$ EVIDÊNCIA $\rightarrow$
    VALIDAÇÃO-CRUZADA $\rightarrow$ CLASSIFICAR, resistente à injeção de logs),
    M4 (Monitoramento Comportamental) e M5 (Quarentena).

    \item \textbf{Inexistência de avaliação empírica de resiliência sobre telemetria
    viva.} As defesas, quando existem, são avaliadas em simulação; as caças, quando
    operam sobre telemetria real \cite{ali2023huntgptintegratingmachinelearningbased},
    não são submetidas a um adversário que ataque o agente. O MAS-Hunt fecha essa
    lacuna ao ser avaliado em laboratório instrumentado com a pilha Elastic,
    confrontando o sistema completo (C1-full) contra a ablação sem endurecimento
    (C1-nohardening), contra uma linha de base de passagem única sem governança
    (C2-naive) e contra as regras padrão do Elastic Security (C3-Elastic), medindo
    tanto desempenho de detecção (Precisão, Revocação, F2) quanto resiliência
    (\textit{Attack Success Rate}) sob injeção adversarial, com intervalos de
    confiança por \textit{bootstrap}.
\end{enumerate}

Em suma, enquanto os trabalhos relacionados respondem isoladamente às perguntas
``como coordenar agentes?'', ``como governar IA?'', ``como atacar agentes?'' e
``como automatizar detecção?'', o MAS-Hunt responde à pergunta integradora que
nenhum deles enfrenta: \textit{como projetar e validar empiricamente um sistema
multiagente de caça a ameaças cuja própria estrutura de governança o torna
resiliente a ataques adversariais dirigidos aos agentes, operando sobre telemetria
de produção?} É essa integração --- e sua verificação experimental --- que
constitui a contribuição original posicionada pelos capítulos seguintes.
